\section{Fixed vocabularies and model checking}\label{sec:fixed-vocabulary}
The lower bound for $G_k$ uses a growing input vocabulary. We now show
that a single binary relation suffices for a superpolynomial composition
lower bound, and that fixing the vocabulary changes the optimal circuit
rate to $2^{\Theta(m/\log(m+2))}$. Throughout this section, $m$ counts
nodes in the \emph{source syntax tree}; only the target may share
subexpressions. Logarithms are to base two.

\begin{theorem}[Fixed-vocabulary succinctness]
\label{thm:fixed-vocabulary:formal}
For every sufficiently large $m$, there is an \FOthree\ query of
syntax-tree size at most $m$ over one binary relation such that every
equivalent relational circuit over finite structures has at least
$2^{\Omega(m/\log(m+2))}$ composition gates.
Conversely, over any fixed finite binary vocabulary $\sigma$, every
\FOthree\ formula of syntax-tree size $m$ with at most two free variables
has an equivalent relational circuit of total size
$2^{O_\sigma(m/\log(m+2))}$, constructible within the same time bound.
The upper-bound translation is correct on all nonempty structures.
\end{theorem}

We first give the constructive upper bound, then transfer the main
lower bound through a single-relation encoding. The final subsection
derives the model-checking consequence.

\subsection{The fixed-vocabulary upper bound}\label{sec:fixed-upper}
Fix a finite binary vocabulary $\sigma$. Since the variable names
are restricted to $x,y,z$, formulas are trees over a fixed finite
ranked alphabet. Consequently, for some constant $a_\sigma>1$, there
are at most $a_\sigma^{t+1}$ syntactically distinct formulas of size
at most $t$. It follows that $d$ distinct formulas
$\psi_1,\ldots,\psi_d$ satisfy
\[
 \sum_{j=1}^d|\psi_j|
 \ge c_\sigma d\log(d+1)
\]
for a constant $c_\sigma>0$.
Indeed, for all sufficiently large $d$, at most $d/2$ formulas have
size below a sufficiently small constant times $\log d$.
At least $d/2$ of the chosen formulas therefore have size
$\Omega_\sigma(\log d)$; reducing $c_\sigma$ handles the remaining
values of $d$.

\begin{proof}[Proof of the upper bound in Theorem~\ref{thm:fixed-vocabulary:formal}]
Use the Boolean-shell translation of
Appendix~\ref{app:classical-translation}, but share the translated
children instead of copying them into each DNF clause.
For a quantified subformula $Qw\,\psi$, its shell treats every atom
and every maximal quantified subformula of $\psi$ as a Boolean input.
Let $d$ be the number of \emph{distinct} quantified children.
Choosing one occurrence of each gives disjoint subtrees of the source,
whose total size is at most $m$. The preceding inequality yields
\[
 d\log(d+1)=O_\sigma(m),
 \qquad d=O_\sigma\!\left(\frac{m}{\log(m+2)}\right).
\]
There are only $O_\sigma(1)$ possible relation and equality atoms.
Thus the shell is a Boolean function of
$r=d+O_\sigma(1)=O_\sigma(m/\log(m+2))$ distinct inputs.
Repeated occurrences of the same child are assigned the same input.

Process quantified subformulas from the inside out. By the
binary-assignment convention of \S\ref{sec:prelim-logic}, each
quantified child has at most two free variables and already has a
binary representation. A unary child has a dummy coordinate, so its
orientation can be changed by converse; a sentence is represented
by $\zero$ or $\one$. Reflexive atoms have constant-size
representations as in \S\ref{sec:FOtwo-translation}.

The truth-table DNF of a shell has at most $2^r$ clauses and $r$
literals per clause. At an existential quantifier, group the literals
of a clause by their free-variable pairs, obtaining
$A(u,v)\land B(u,w)\land D(w,v)$. Eliminating $w$ gives
$A\cap(B;D)$, with empty groups represented by $\one$.
Thus each clause introduces at most one composition.
Negated literals and universal quantifiers use complement.
Each child circuit is computed once and reused in every clause:
only $O((r+1)2^r)$ new gates, including at most $2^r$ compositions,
are introduced at this shell.

There are at most $m$ quantified occurrences. The unquantified root
shell, whose free variables are contained in a pair, is translated
directly by Boolean operations. All atomic representations together
cost $O_\sigma(1)$ gates if shared. Hence both the total size and the
composition count are bounded by
\[
 |C_\varphi|
 \le \operatorname{poly}_\sigma(m)\,
       2^{O_\sigma(m/\log(m+2))}
 =2^{O_\sigma(m/\log(m+2))}.
\]
The translation is correct on all nonempty structures by the same
quantifier-elimination identity as the classical translation.
Syntactic comparison of children, enumeration of shell truth
assignments, and construction of the circuit take at most a polynomial
factor in $m$ times this bound, which is absorbed into the exponent.
\end{proof}

The counting step uses disjoint subtrees of the source. It therefore
does not give the same upper bound when $m$ instead measures a
DAG-represented source formula. Nor does the construction guarantee a
subexponential-size \CoR\ term: unfolding the shared children can
increase its size.

\subsection{The lower bound over one binary relation}\label{sec:single-relation}
We encode the input relations of $G_k$ in a single relation on copies
of the original domain. The following block calculation preserves
the composition measure even with unrestricted complement and sharing.

\begin{lemma}[Block simulation]\label{lem:block-simulation}
Suppose an encoding maps each structure with domain $V$ to a structure
on $V\times L$, where $|L|=q$, with one input relation $E$.
Assume that each block
$E_{ab}(u,v)\coloneqq E((u,a),(v,b))$ is an original input relation or one of
$\zero,\one,\id$, with the choice depending only on $a,b\in L$.
Every relational circuit $C(E)$ can then be simulated over the
original signature so as to compute any designated output block,
using at most $q^3\gcost C$ composition gates.
\end{lemma}
\begin{proof}
For every gate $X$, compute all $q^2$ relations
$X_{ab}(u,v)\coloneqq X((u,a),(v,b))$.
The blocks of $\zero$ and $\one$ are respectively $\zero$ and $\one$;
the identity has block $\id$ when $a=b$ and $\zero$ otherwise.
Boolean operations act blockwise, including complement relative to
$V^2$, and converse exchanges the layer indices:
\[
 (\overline X)_{ab}=\overline{X_{ab}},
 \qquad
 (X^\smile)_{ab}=(X_{ba})^\smile.
\]
Composition satisfies
\[
 (X;Y)_{ab}=\bigcup_{c\in L}(X_{ac};Y_{cb}).
\]
Thus each composition gate introduces at most $q^3$ compositions
over all output blocks; the other operations introduce none.
The blocks of a shared gate are computed once and reused.
Induction over the circuit proves correctness.
\end{proof}

\paragraph{The encoding.}
List the $2k$ input symbols of $G_k$ as
$R_1,\ldots,R_{2k}$, with $R_i\coloneqq P_i$ and $R_{k+i}\coloneqq Q_i$.
Put
\[
 \ell\coloneqq\lceil\log(2k)\rceil,\qquad
 L\coloneqq\{0,1\}^{\le\ell},\qquad
 q\coloneqq|L|=2^{\ell+1}-1<8k,
\]
and choose distinct leaves $w_1,\ldots,w_{2k}\in\{0,1\}^{\ell}$.
For a structure $\mathcal M$ on $V$, the encoded structure
$\operatorname{Enc}(\mathcal M)$ has domain $V\times L$.
Its sole relation $E$ consists of the following edges:
\[
\begin{aligned}
 &((u,t),(u,tb))
       &&(u\in V,\ |t|<\ell,\ b\in\{0,1\}),\\
 &((u,t),(u,t))
       &&(u\in V,\ t\ne\epsilon,\ t\text{ ends in }1),\\
 &((u,w_j),(v,\epsilon))
       &&((u,v)\in R_j^{\mathcal M},\ 1\le j\le2k).
\end{aligned}
\]
There are no other edges. Each $u$ has a binary routing tree; loops
mark its $1$-children, and the edges from leaf $w_j$ back to the root
layer store $R_j$.
Every layer-to-layer block is $\zero$, $\id$, or some $R_j$, so
Lemma~\ref{lem:block-simulation} applies.
The encoded graphs are directed and may have loops.

\paragraph{Three-variable decoders.}
Let
\[
 N(u,v)\coloneqq E(u,v)\land u\ne v,\qquad
 L_0(v)\coloneqq\neg E(v,v),\qquad L_1(v)\coloneqq E(v,v).
\]
For distinct variable names $u,v\in\{x,y,z\}$, define
\[
 \theta_\epsilon(u,v)\coloneqq N(u,v),\qquad
 \theta_{bs}(u,v)\coloneqq
 \exists w\bigl(N(u,w)\land L_b(w)\land\theta_s(w,v)\bigr),
\]
where $w$ is the third variable name.
Inside $\theta_s(w,v)$, the old name $u$ may be rebound; all such
bindings are local to that subformula.
For every word $t$, expanding $\theta_t$ gives a genuine
\FOthree\ formula of size $\Theta(|t|+1)$.

Starting at $(u,\epsilon)$, the first $\ell$ non-loop steps of
$\theta_{w_j}$ uniquely follow the word $w_j$ through the routing tree.
Its final non-loop edge returns to the root layer and reads $R_j$.
Consequently,
\[
 \operatorname{Enc}(\mathcal M)\models
 \theta_{w_j}((u,\epsilon),(v,\epsilon))
 \quad\Longleftrightarrow\quad
 \mathcal M\models R_j(u,v).
\]
Furthermore, each selected decoder $\theta_{w_j}$ starting at the
root layer can hold only when its second argument is also in that layer.
The inequalities in $N$ exclude the marking loops from both the
routing steps and the final data step; data edges can still encode
loops of the original relations, since their two layer indices differ.

\begin{proof}[Proof of the lower bound in Theorem~\ref{thm:fixed-vocabulary:formal}]
Replace each relation-atom occurrence $R_j(u,v)$ in $G_k$ by
$\theta_{w_j}(u,v)$, obtaining $\widehat G_k$ over the single symbol $E$.
Each decoder uses only three variable names. Any rebinding of the
third name is local to that atom replacement, so it does not affect
another conjunct or disjunct of $G_k$.
There are $4k$ relation-atom occurrences, hence
\[
 m_k\coloneqq|\widehat G_k|=\Theta(k\log(k+1)).
\]
When $x,y$ lie in the root layer, every successful decoded atom on
$(x,z)$ forces $z$ into the root layer as well. Since $k\ge1$,
\[
 \operatorname{Enc}(\mathcal M)\models
 \widehat G_k((u,\epsilon),(v,\epsilon))
 \quad\Longleftrightarrow\quad
 \mathcal M\models G_k(u,v).
\]

Let $C$ compute $\widehat G_k$ on all finite $E$-structures.
Lemma~\ref{lem:block-simulation} produces a circuit $C^\downarrow$
over the original signature computing its root--root output block,
with $\gcost{C^\downarrow}\le q^3\gcost C$.
The displayed equivalence gives $C^\downarrow\eqfin G_k$, so
Theorem~\ref{thm:main} yields
\[
 2^{k-1}\le\gcost{C^\downarrow}\le q^3\gcost C,
 \qquad
 \gcost C\ge\frac{2^{k-1}}{q^3}
       >\frac{2^{k-1}}{(8k)^3}.
\]
Since $k=\Theta(m_k/\log(m_k+2))$, this is
$2^{\Omega(m_k/\log(m_k+2))}$.
For every sufficiently large size bound $m$, choose
$k=\lfloor c m/\log(m+2)\rfloor$ with a sufficiently small absolute
constant $c>0$. Then $m_k\le m$ and the same estimate gives the
stated $2^{\Omega(m/\log(m+2))}$ lower bound.
No well-formedness test for the encoding is needed: correctness on
all finite $E$-structures already includes the encoded ones.
\end{proof}

For any fixed vocabulary containing at least one binary symbol, the
same lower bound follows by fixing all other input relations to
$\zero$. Together with the upper bound, this determines both the
worst-case composition count and the worst-case total circuit size
at the scale $2^{\Theta(m/\log(m+2))}$.
The composition lower bound also holds for terms, but does not
determine their optimal general upper bound.

\subsection{Model checking on directed graphs}\label{sec:fixed-model-checking}
\begin{corollary}\label{cor:fixed-model-checking}
Fix a finite binary vocabulary $\sigma$ and $\eta>0$.
A tree-represented $\mathrm{FO}^3[\sigma]$ sentence of size $m$ can be
evaluated on an $n$-element structure in
\[
 2^{O_\sigma(m/\log(m+2))}\,n^{\omega+\eta}
\]
time, including compilation. In particular, this holds for directed
graphs with one input edge relation.
\end{corollary}
\begin{proof}
Construct the circuit of Theorem~\ref{thm:fixed-vocabulary:formal} and evaluate
each gate once in topological order.
A composition is a Boolean matrix product, computable in
$O(n^{\omega+\eta})$ time for any fixed $\eta>0$; all other gates take
$O(n^2)$ time. By the assignment semantics of \S\ref{sec:prelim-logic}, a
sentence denotes either $\one$ or $\zero$, so any output entry
gives its truth value.
The circuit construction time and input-reading time are absorbed
in the displayed bound.
\end{proof}

\paragraph{Conditional tightness.}
For fixed vocabularies, this improves the $2^{O(m)}$ formula-size
factor in~\citep[Theorem~10]{nakamura-uezato2027}.
Under the low-dimensional 3-OV hypothesis, their Theorem~72
(Appendix~T) rules out $2^{o(m/\log(m+2))}n^{3-\delta}$ time
for every fixed $\delta>0$, even for tree-represented
existential-positive sentences over one binary relation.
Their bound counts edges $M$ rather than vertices.
Retaining at most three isolated vertices preserves \FOthree\
sentences and leaves at most $2M+3$ vertices
\citep[Proposition~6]{nakamura-uezato2027}, giving the stated vertex
bound. Building an adjacency matrix takes $O((M+1)^2)$ time,
absorbed after reducing $\delta$ to lie in $(0,1)$.
Thus $m/\log m$ is the conditionally optimal exponent scale for
truly subcubic evaluation, without a claim of optimality for $\omega$.
